%% =========================================================================
%% When Stronger Text Evidence Hurts:
%% Fallback-Safe Conditioning of Discrete Preference Diffusion Kernels
%%
%% DRAFT v2 (repositioned), created 2026-07-06.
%%   - Supersedes paper/main.tex (frozen as backup; do not edit it).
%%   - Claim tier: Family D baseline per
%%     docs/reports/2026-07-04-family-d-claim-freeze-cn.md, PLUS the
%%     Gate-2 amendment: one closed-gate ML1M miss plus one barely-open
%%     ATG miss remain open implementation questions. ML1M exceeds its
%%     measured host floor; ATG exceeds only a provenance-limited observed
%%     three-run spread because the run manifests are absent. E1 found a
%%     step-0 optimizer-ownership mismatch, so end-to-end empirical
%%     reduction wording stays paused.
%%   - Repositioning (spec 2026-07-06 section 5.1): intro paragraph and a
%%     fifth related-work axis situate the paper against the discrete-
%%     diffusion generative-recommendation wave; claim sentences unchanged.
%%
%% HOW TO BUILD
%%   1. Download the AAAI-27 author kit and place aaai27.sty (+ aaai27.bst)
%%      next to this file. The draft falls back to a plain two-column
%%      approximation when the kit is absent, so it compiles either way.
%%   2. pdflatex main_v2 && bibtex main_v2 && pdflatex main_v2 && pdflatex main_v2
%%
%% PLACEHOLDER POLICY (unchanged): every experiment-dependent value still
%%   pending is marked \pending{...} or \pnum. Placeholders may ONLY be
%%   replaced from dated artifacts under docs/reports/data (cite the
%%   artifact in a % comment when filling). Owner of backfill: CLOSE-07.
%% =========================================================================
\documentclass[letterpaper]{article}

\IfFileExists{aaai27.sty}{%
  \usepackage{aaai27}%
  \usepackage{times}%
  \usepackage{helvet}%
  \usepackage{courier}%
}{%
  % ---- fallback preamble (kit absent); harmless to delete once kit is in
  \usepackage[margin=0.75in]{geometry}%
  \usepackage{multicol}%
  \usepackage{times}%
}

\usepackage{graphicx}
\usepackage{amsmath,amssymb,amsthm}
\usepackage{booktabs}
\usepackage{xcolor}
\usepackage{url}

% ---------------- placeholder macros (must all be gone before submission)
\newcommand{\pending}[1]{{\color{red}\textbf{[PENDING: #1]}}}
\newcommand{\pnum}{{\color{red}$\boxdot$}}         % pending numeric cell
\newcommand{\pfig}[1]{{\color{red}\fbox{\parbox{0.9\linewidth}{\centering\vspace{2em}FIGURE PENDING: #1\vspace{2em}}}}}

% ---------------- notation
\newcommand{\pcore}{p_{\mathrm{core}}}
\newcommand{\pcorer}{\tilde{p}_{\mathrm{core}}}
\newcommand{\ch}{c_{h}}
\newcommand{\pu}{p_{u}}
\newcommand{\ut}{\tilde{u}}
\newcommand{\Uds}{U_{\mathrm{ds}}}
\newcommand{\phids}{\phi(\Uds)}
\newcommand{\TV}{\mathrm{TV}}
\newtheorem{proposition}{Proposition}
\newtheorem{lemma}{Lemma}
\newtheorem{remark}{Remark}
\newtheorem{corollary}{Corollary}

\title{When Stronger Text Evidence Hurts: Fallback-Safe Conditioning of
Discrete Preference Diffusion Kernels}
% Frozen alt titles (4.1 candidates; CLOSE-06 owns the final pick):
%   Fallback-Safe Evidence Conditioning for Discrete Preference Diffusion
%   Reliability-Gated Proposal Kernels for Discrete Preference Diffusion
\author{Anonymous submission}

\begin{document}
\maketitle

\begin{abstract}
Discrete preference diffusion recommenders such as PreferGrow corrupt a
user's preferred item by fading it toward a \emph{proposal distribution}
and learn to grow the preference back. Personalization in this family
acts entirely on the \emph{score} side: the transition kernel is shared
by all users. We personalize the kernel, and in doing so surface a
counterintuitive fact: because the proposal is a \emph{corruption}
distribution, text evidence can be too good to use. On four benchmarks,
a train-only measure of how well text similarity predicts the next item
inversely rank-orders the outcome of text-tilting the kernel in all four
cases: the tilt helps where text defines informative hard negatives and
fails where text predicts future positives; corruption then becomes
false-negative pollution. We operationalize this with a two-factor,
history-only gate (a per-user coherence signal calibrated against a
length-matched random-history null, damped by a frozen dataset-level
utility factor) that interpolates between the host's learned proposal
and a frozen text anchor. The design is \emph{fallback-safe} at the
kernel level: the conditioned chain retains Markovianity, detailed
balance, and the analytic reversal of the host kernel; it reduces
\emph{exactly} to the host when the gate closes. The bounded-deviation
claim concerns only the proposal and one forward transition row:
$\TV(\pu,\pcore)\le g$ and
$\TV(P_{\sigma,\pu}(i,\cdot),P_{\sigma,\pcore}(i,\cdot))\le
(1-\alpha_\sigma)g$; it is not a bound on end-to-end recommendation
loss. Controls show that the mechanism requires an \emph{aligned}
per-user signal. A pre-registered four-dataset validation returned a
mixed single-run observation: fully open Steam kept a small positive
gain, closed-gate Beauty held its parity band, while the closed-gate
ML1M miss exceeded its measured host noise floor. The barely-open ATG
($\phids{=}0.117$) miss was $4.37\times$ a provenance-limited observed
three-run spread, but all three run manifests are absent, so we do not
call that spread a fully manifested noise floor. E1 also found a step-0
optimizer-ownership mismatch between the production paths. We report
the two misses as open implementation questions rather than evidence
in either direction about gate efficacy. The cross-dataset inversion rests on four points;
where text utility is high, the design closes the kernel toward the
host rather than promising an empirical gain.
% Frozen Family D main + boundary sentences implemented above; freeze memo:
% docs/reports/2026-07-04-family-d-claim-freeze-cn.md (upgrade only via
% documented amendment at CLOSE-05).
\end{abstract}

\section{Introduction}

Diffusion-based recommenders replace single-step prediction with an
iterative corrupt-and-reconstruct process. In the discrete
instantiation (preference fading and growing over the item
vocabulary~\cite{prefergrow}, built on structured discrete kernels and
score-entropy training~\cite{d3pm,sedd}), corruption acts directly on
the object being ranked: the forward kernel fades the preferred item
into a \emph{proposal distribution} over alternatives, and the reverse
process grows the preference back under a learned score model.

The proposal distribution is the central object of this family. It
encodes ``what replaces a faded preference,'' and through that choice
it determines the stationary distribution of the forward chain, the
reference measure of its detailed balance, the weights of the
score-entropy objective, and the analytic reverse ratio. Yet in
existing discrete preference diffusion the proposal is \emph{global}:
one kernel shared by every user, with personalization entering only
through the score network and inference-time guidance against a
``non-preference user'' branch~\cite{prefergrow}. Meanwhile, item
text, the most abundant external evidence about items, is routinely
exploited on the representation side of sequential
recommenders~\cite{recformer,unisrec}, and continuous-diffusion
recommenders condition denoisers or guidance
signals~\cite{dreamrec,preferdiff,diffurec}. The corruption kernel
itself has remained unconditioned.

Making corruption \emph{adaptive} is now an active program in
generative recommendation: parallel semantic-ID decoding under masked
diffusion~\cite{llada-rec}, confidence-ordered any-order
decoding~\cite{gr-diffgrm}, and position-adaptive masking schedules
motivated by the observation that uniform masking is
harmful~\cite{gr-admasked}, with theoretical analysis beginning to
follow~\cite{gr-mdtheory} on top of the discrete-diffusion
language-model program~\cite{dllm-llada}. What this wave adapts to,
however, is model-\emph{internal} signal: confidence and position
heuristics. Two ingredients are missing for \emph{external} evidence:
a train-computable decision rule for whether evidence should be
allowed to shape corruption at all, and a guarantee for the shaped
process when the evidence is misleading. This paper is a
pre-registered reliability study of exactly those two ingredients, in
the simplest kernel family where they can be posed cleanly.

Intuition suggests a monotone rule: the better the text evidence, the
harder the kernel should lean on it. \emph{Our central empirical
finding is that this intuition is backwards.} The proposal is a
corruption distribution, so the items it favors are treated as ``where
preference goes when it fades,'' that is, as negatives. If text
similarity is informative enough to rank the user's \emph{actual next
item} highly, a text-tilted proposal concentrates corruption mass on
likely future positives and manufactures false negatives in proportion
to the quality of the evidence, a failure mode the negative-sampling
literature pays dearly to avoid~\cite{srns,debiased}. Measuring
evidence quality with a train-only statistic $\Uds$ (how well text
similarity discriminates the next item from sampled negatives), we
find that $\Uds$ \emph{inversely} rank-orders the outcome of
text-tilting the kernel on all four of our benchmarks
(Sec.~\ref{sec:inversion}): the largest gain occurs on the dataset
with the \emph{weakest} next-item text signal, and a severe loss on
the dataset with the \emph{strongest}.

The inversion dictates what a reliability signal for kernel
conditioning must do, and motivates our design, \emph{fallback-safe
evidence conditioning}. A low-capacity gate
\begin{equation*}
  g(\ut,\Uds) \;=\; g_{\max}\;\phids\;\mathrm{clip}(\ut,0,1)
\end{equation*}
combines two history-only factors: a per-user coherence signal $\ut$,
calibrated against a length-matched \emph{random-history null} of the
evidence space (raw coherence has a high, dataset-dependent floor that
otherwise couples the kernel's operating point to embedding geometry);
and a frozen dataset-level utility factor $\phids$ that \emph{closes}
the gate as $\Uds$ enters the false-negative regime. The gated kernel
interpolates between the host's \emph{learned} global proposal and a
frozen text anchor and leaves the host's pseudo-item mass untouched.
Conditional on the fixed history-derived signals, all structural
properties of the host chain survive; when the gate closes, the
proposal and its derived kernel objects reduce exactly to the host.
For nonzero gates, Lemma~\ref{lem:tv} bounds the proposal and each
one-step forward transition row in total variation; it does not bound
the trained score model or recommendation metrics. The theory is
deliberately elementary (rank-one kernel algebra), but it buys a
property that score-side or loss-side side-information injection
cannot even express: a provable reduction to the host when evidence is
uninformative \emph{or misleading}. Neither ingredient is tied to the
host: $\Uds$ uses only the evidence bank and training transitions, and
the gate and its guarantees are stated for rank-one absorbing kernels
rather than for one architecture. A pre-registered four-dataset
validation and mechanism controls then test the design under a frozen
protocol. The single-run readout contains one closed-gate parity hit
(Beauty), one closed-gate miss (ML1M), one barely-open miss (ATG,
$\phids{=}0.117$), and a fully open directional gain that misses its
reference magnitude (Steam). We keep the two misses as open
implementation questions and do not treat them as an end-to-end
demonstration of exact reduction.

Our contributions:
\begin{itemize}
  \item \textbf{The utility-inversion finding.} A train-only text
  utility statistic inversely rank-orders the outcome of text-tilting a
  discrete preference-diffusion kernel on all four benchmarks. Under
  exchangeability, one specified order among $4!$ equally likely label
  orders gives the descriptive fraction $1/4!=1/24$; this is not a
  confirmatory p-value. We give a mechanism-level account: corruption
  proposals reward hard negatives and punish predicted positives
  (Sec.~\ref{sec:inversion}). We
  further propose $\Uds$ as a train-only \emph{preflight} statistic
  for any scheme that mines negatives from content similarity; our
  kernel results are its first evidence.
  % FOLLOWUP-09 did not confirm the inversion at user level, so the
  % preflight proposal remains supported at the cross-dataset level
  % only in this draft.
  \item \textbf{Kernel-side personalization with fallback safety.} The
  first discrete preference-diffusion recommender whose transition
  kernel is conditioned per user by an explicit, history-only evidence
  signal; conditional Markovianity, the stationary proposal, detailed
  balance, and the analytic reverse ratio all carry over
  (Prop.~\ref{prop:validity}), with exact reduction at $g=0$
  (Prop.~\ref{prop:reduction}) and nonzero-gate bounds limited to
  proposal TV and one-step forward transition-row TV
  (Lemma~\ref{lem:tv}), not end-to-end recommendation performance.
  \item \textbf{Calibration methodology.} A length-matched
  random-history null for user-level coherence---raw and even
  null-calibrated coherence mismeasure utility, which is why the
  dataset-level factor is estimated from utility directly and frozen
  before any retraining (pre-registered thresholds,
  Appendix~\ref{app:uds}).
  \item \textbf{Mechanism evidence.} Controls (reliability shuffle,
  global proposal, ungated anchor) show that aligned per-user gating is
  load-bearing. First-generation corruption diagnostics are reported as
  limited observations rather than evidence for a final-v2
  corruption-response claim
  (Sec.~\ref{sec:mechanism}).
\end{itemize}

\section{Preliminaries: Discrete Preference Diffusion}
\label{sec:prelim}

Let $\Omega=\{1,\dots,M\}$ be the item vocabulary; following
PreferGrow~\cite{prefergrow}, an optional non-preference pseudo-item is
appended as the last coordinate, giving $|\Omega|=M{+}1$ states. We use
the row-vector convention: a distribution $\mu$ evolves as $\mu P$.

Given a proposal distribution $p$ on $\Omega$, define the rank-one
kernel $\Pi_p(i,j)=p(j)$ and the generator $Q_p=\Pi_p-I$. With noise
scale $\sigma\ge 0$ and $\alpha_\sigma=e^{-\sigma}$, the forward
transition is
\begin{equation}
  P_{\sigma,p} \;=\; e^{\sigma Q_p} \;=\; \alpha_\sigma I +
  (1-\alpha_\sigma)\,\Pi_p ,
  \label{eq:forward}
\end{equation}
so a preferred item survives with probability $\alpha_\sigma$ and is
otherwise resampled from $p$: preference \emph{fades toward the
proposal}. The score model $s_\Theta(x_t,t,u)$ (with $u$ a SASRec-style
history encoding) is trained by denoising score entropy~\cite{sedd},
and the reverse process uses the analytic stationary-reversal ratio in
which the proposal appears in closed form (Appendix~\ref{app:proofs}).
The host learns its proposal: $\pcore=\mathrm{softmax}(p_1)$ with $p_1$
a free parameter vector, trained jointly. Everything below treats the
host score model, objective, and schedule as fixed; we modify
\emph{only} the proposal.

\section{Method: Fallback-Safe Gated Proposals}
\label{sec:method}

Figure~\ref{fig:method} overviews the construction; the pieces are
specified below.

\begin{figure}[t]
\centering
\pfig{Figure 1 (method overview): draw per the Figure-1 prompt in
docs/reports/2026-07-08-figure-prompts.md, then replace this box with
an includegraphics call.}
\caption{Overview of fallback-safe evidence conditioning. A frozen
text bank defines a per-history content anchor $\ch$; a two-factor,
history-only gate $g=g_{\max}\,\phids\,\mathrm{clip}(\ut,0,1)$
interpolates the corruption proposal between the host's learned
$\pcore$ and $\ch$. When the gate closes ($g{=}0$) the kernel is
exactly the host's (Prop.~\ref{prop:reduction}); otherwise,
Lemma~\ref{lem:tv} bounds proposal TV by $g$ and each one-step forward
transition-row TV by $(1-\alpha_\sigma)g$. These are kernel-level
bounds, not bounds on NDCG or complete-trajectory behavior.}
\label{fig:method}
\end{figure}

\subsection{Evidence anchor}
Each dataset has a frozen text bank: one record per item (title, brand,
category, description where present), embedded once by an off-the-shelf
text encoder and never fine-tuned. For a user history
$h=(i_1,\dots,i_L)$ let $\bar e_h$ be the normalized mean of its item
embeddings and define the \emph{content anchor}
\begin{equation}
  \ch(i) \;\propto\; \exp\!\big(\langle \bar e_h, e_i\rangle/\tau_0\big),
  \label{eq:anchor}
\end{equation}
a distribution over real items, down-weighted for items with missing
text. The anchor proposes \emph{content-plausible substitutes}; it
never replaces the host's score representation.

\subsection{User factor: null-calibrated coherence}
\label{sec:null}
The raw coherence statistic
$A_h=\frac{1}{L}\sum_{\ell}\tfrac{1}{2}\big(\langle e_{i_\ell},\bar
e_h\rangle+1\big)$ has a high, dataset-dependent floor: on our
production banks a \emph{random} pseudo-history already scores
$0.931$ on ML1M, $0.942$ on both Steam and Beauty
% artifact: docs/reports/data/2026-07-02-gate0/gate0_failure_component_summary.csv
because the statistic inherits the global geometry of the embedding
space. Consuming $A_h$ in absolute terms couples the kernel's operating
point to that geometry. We therefore calibrate against the
evidence-space null: for each dataset and history length $L$ we draw
random pseudo-histories from the frozen bank, record
$\mu_{\mathrm{null}}(L),\sigma_{\mathrm{null}}(L)$, and define
\begin{equation}
  \ut \;=\;
  \frac{A_h-\mu_{\mathrm{null}}(L)}{k\,\sigma_{\mathrm{null}}(L)} .
  \label{eq:utilde}
\end{equation}
$\ut$ is history-only (no target, no diffusion state, no model score)
and fixed along each forward/reverse trajectory. A history at or below
its null contributes zero gate.

Null calibration is necessary but \emph{not sufficient}: calibrated
coherence measures whether behavior clusters in text space, not whether
tilting corruption toward that cluster helps. On our production banks
the median $\ut$ ranks ML1M ($1.43$) $>$ Beauty ($0.80$) $>$ ATG
($0.72$) $>$ Steam ($0.11$)
% artifact: docs/reports/data/2026-07-02-gate0/gate0_u_tilde_summary.csv
; this is nearly the \emph{opposite} of where text tilt actually helps
(Sec.~\ref{sec:inversion}). This mismatch is exactly why a second,
utility-facing factor is required.

\subsection{Dataset factor: frozen text utility}
\label{sec:uds}
We measure directly on the training split what the gate actually
needs to know: does text similarity predict \emph{positives}? For sampled
train transitions $(h,\text{next})$, $\Uds$ is the probability that the
anchor similarity ranks the true next item above a popularity-sampled
negative (an AUC; estimator frozen in Appendix~\ref{app:uds}). Like the
host's popularity prior, $\Uds$ is a train-only dataset constant; no
validation or test information enters. The gate's dataset factor is the
frozen hinge
\begin{equation}
  \phids \;=\; \mathrm{clip}\!\big((0.70-\Uds)/0.10,\;0,\;1\big),
  \label{eq:phi}
\end{equation}
fully open when $\Uds\le 0.60$ (text defines hard negatives but does
not predict positives) and fully closed when $\Uds\ge 0.70$ (tilting
would concentrate corruption on likely positives). The thresholds were
fixed \emph{before} the production statistics were computed and were
not adjusted afterwards (pre-registration note in
Appendix~\ref{app:uds}).

\subsection{The gated kernel}
Let $\pcore=\mathrm{softmax}(p_1)$ be the host's learned proposal
(including the pseudo-item coordinate $\star$) and $\pcorer$ its
restriction to real items, renormalized. With
$g=g_{\max}\,\phids\,\mathrm{clip}(\ut,0,1)$,
\begin{equation}
\begin{aligned}
  \pu \;=\;& \pcore(\star)\, e_\star \\
  &+ \big(1-\pcore(\star)\big)\Big[(1-g)\,\pcorer + g\,\ch\Big],
\end{aligned}
\label{eq:pu}
\end{equation}
and $\pu=(1-g)\pcore+g\,\ch$ when no pseudo-item is used. Three
deliberate restrictions: (i) the \emph{pseudo-item mass is exactly the
host's} and the gate never touches it; (ii) $g$ is the \emph{only}
dial, with no evidence-dependent temperature or popularity mixing; (iii) no
new learnable modules: the only trained parameter in the proposal path
is the host's own $p_1$, shared with the fallback branch. $\pu$
replaces the global proposal in \emph{every} kernel-level object:
forward corruption \eqref{eq:forward}, score-entropy weights, and the
analytic reverse ratio, not just negative sampling.

\subsection{Theory}
\label{sec:theory}
Proofs are in Appendix~\ref{app:proofs}. Fix a user history $h$, hence
$\ut$, $\Uds$, and $\pu$, for the whole trajectory.

\begin{proposition}[Kernel validity]
\label{prop:validity}
With $\pu$ in place of the global proposal, the forward corruption has
the closed form
$q_t(x_t{=}x \mid x_0,u)=\bar\alpha_t\mathbf{1}\{x{=}x_0\}
+(1-\bar\alpha_t)\pu(x)$; conditional on the fixed gate inputs the
chain is Markov; its stationary and limiting distribution is $\pu$; it
satisfies detailed balance w.r.t.\ $\pu$; and the analytic
stationary-reversal ratio is the host's formula with the proposal term
replaced by $\pu$.
\end{proposition}
\begin{proposition}[Exact reduction / fallback safety]
\label{prop:reduction}
If $g=0$ then $\pu=\pcore$, hence $Q_{\pu}=Q_{\pcore}$ and every
derived object (forward corruption, stationary distribution,
score-entropy weights, analytic reverse ratio) coincides exactly with
the host kernel's.
\end{proposition}
\begin{lemma}[Bounded deviation]
\label{lem:tv}
$\TV(\pu,\pcore)\le g$, and each row of the forward transition
satisfies
$\TV\big(P_{\sigma,\pu}(i,\cdot),P_{\sigma,\pcore}(i,\cdot)\big)\le
(1-\alpha_\sigma)\,g$.
\end{lemma}
\begin{remark}[Why the kernel is the right injection site]
\label{rem:site}
Encoder- or loss-level injection of the same signals admits no analogue
of Prop.~\ref{prop:reduction}: there is no sense in which ``the encoder
with text features at zero reliability'' \emph{is} the host model. The
reduction guarantee is expressible only where conditioning enters the
transition law itself. Our injection-site comparison is therefore an
analysis of regimes, not the basis of the safety claim.
\end{remark}
The guarantees hold under the stated scope: gate inputs history-only
and frozen per trajectory, self-resampling corruption, and the
pseudo-item coordinate present in every kernel object
(Appendix~\ref{app:boundary}).

\section{The Utility-Inversion Analysis}
\label{sec:inversion}

\begin{figure}[t]
\centering
\pfig{Figure 2 (inversion mechanism): draw per the Figure-2 prompt in
docs/reports/2026-07-08-figure-prompts.md, then replace this box with
an includegraphics call.}
\caption{Why stronger text evidence hurts. Left (low $\Uds$): text
similarity structures the item space but does not predict the next
item, so tilting corruption toward the anchor supplies informative
hard negatives. Right (high $\Uds$): text similarity ranks the true
next item highly, so the same tilt lands corruption on likely future
positives and manufactures false negatives
(Prop.~\ref{prop:collision}). The frozen hinge $\phids$ closes the
gate in the right-hand regime.}
\label{fig:inversion}
\end{figure}

Table~\ref{tab:inversion} is the paper's central empirical exhibit;
Figure~\ref{fig:inversion} sketches the mechanism behind it. For
each dataset we report the frozen production statistic $\Uds$, the
resulting gate factor $\phids$, and the outcome of text-tilting the
kernel \emph{without} a dataset factor: a first-generation,
user-factor-only instantiation of kernel gating trained before $\phids$
existed (parameterization differences in Appendix~\ref{app:gen1};
these numbers motivate the dataset factor and are not a benchmark of
the final system).

\begin{table}[t]
\centering
\small
\begin{tabular}{lccr}
\toprule
Dataset & $\Uds$ & $\phids$ & $\Delta$NDCG@10 (no $\phi$) \\
\midrule
Steam  & 0.570 & 1.000 & $+0.0064$ \\
ATG    & 0.688 & 0.117 & $+0.0019$ \\
Beauty & 0.712 & 0.000 & $-0.0006$ \\
ML1M   & 0.754 & 0.000 & $-0.0622$ \\
\bottomrule
\end{tabular}
\caption{Train-only text utility $\Uds$ (popularity-negative AUC)
inversely rank-orders the test outcome of user-gated kernel tilting
relative to the host, $4/4$. If the four dataset labels are
exchangeable and all $4!$ orders are treated as equally likely, one
specified order has the descriptive fraction $1/24$; this is not a
confirmatory p-value. The frozen hinge $\phids$ closes the gate
precisely where tilting hurt.}
% artifacts: docs/reports/data/2026-07-02-gate0/gate0_text_utility_summary.csv
%            docs/reports/data/2026-07-01-text-side-vs-core-main-table-remote-v77.csv
\label{tab:inversion}
\end{table}

Three readings. \emph{First}, the ordering is perfectly inverted: the
only clear gain (Steam, $+0.0064$) occurs at the \emph{lowest} text
utility, and the severe loss (ML1M, $-0.0622$) at the \emph{highest}.
With four datasets this is suggestive rather than conclusive. Under
exchangeability, one specified order among $4!$ equally likely label
orders gives the descriptive calculation $1/24$; no confirmatory
inference is attached to that number. The ordering is train-computable,
pre-registered into $\phids$, and mechanistically expected: a proposal
is a corruption law, so evidence that predicts positives converts tilt
into false-negative pollution, while evidence that merely structures
the item space supplies informative hard negatives. \emph{Second}, the
inversion is invisible to coherence: ML1M users are the \emph{most}
coherent relative to their null (median $\ut=1.43$), which is why
coherence-only gating cannot save it and why $\phids$ must come from
utility, not coherence. \emph{Third}, popularity does not explain the
signal: $\Uds$ uses popularity-matched negatives, and the uniform-
negative variant gives the same ordering
($0.799/0.692/0.717/0.565$ for ML1M/ATG/Beauty/Steam).
% artifact: gate0_text_utility_summary.csv (u_ds_uniform column)

The false-negative channel can be stated exactly rather than
narratively:
\begin{proposition}[Target collision grows with anchor utility]
\label{prop:collision}
Fix a user with history $h$ and realized next item $y$. Under the
gated kernel, the probability that forward corruption relocates a
faded preference onto $y$ itself is
$\Pr[x_t{=}y \mid x_0{\ne}y] = (1-\bar\alpha_t)\,\pu(y)$ with
\[
\pu(y) \;=\; \big(1-\pcore(\star)\big)\big[(1-g)\,\pcorer(y) +
g\,\ch(y)\big],
\]
affine and strictly increasing in $g$ whenever $\ch(y)>\pcorer(y)$;
averaged over users the slope is
$\mathbb{E}_h\!\big[\ch(y_h)-\pcorer(y_h)\big]$, the mass-form analogue
of the rank statistic $\Uds$ (proof in Appendix~\ref{app:proofs}).
\end{proposition}
Opening the gate on a high-utility bank therefore linearly increases
the rate at which training treats the user's actual next item as
faded-to noise: the increase enters the corruption samples, the
score-entropy weights, and the reverse reference alike. Collision mass is a channel, not by
itself a proof of end-metric harm (the score model can partially
compensate); Table~\ref{tab:inversion} is the empirical arbitration.

\paragraph{Hinge robustness.}
Sweeping the hinge over $U_{hi}\in[0.60,0.82]$ and width
$\in\{0.05,0.10,0.15,0.20\}$ ($92$ grid points), $65\%$ of the grid
preserves both frozen qualitative conclusions (ML1M closed at
$\phi\le0.2$, Steam open at $\phi\ge0.5$); the pre-registered point
$(0.70,0.10)$ lies inside this plateau and admits \emph{all} of the
observed first-generation gain ($0.0066$ summed
$\phi$-weighted) and \emph{none} of the observed harm. Leave-one-out,
the inversion holds in all four cases ($\rho=-1$; each three-point
test is individually weak, exact $p=1/6$).
% artifact: docs/reports/data/2026-07-04-phi-sensitivity-grid.csv
%           docs/reports/data/2026-07-04-phi-sensitivity-summary.md

\paragraph{User-level readout.}
We also ran the within-dataset per-user follow-up on sampled
train-user transitions, regenerating the first-generation readout from
frozen best checkpoints when cached per-user dumps were unavailable.
The inverse utility-harm ordering did \emph{not} reproduce at user
level: Spearman correlations between per-user train-only text utility
and per-user $\Delta\mathrm{NDCG}@10$ were weak and non-negative on
both datasets ($\rho=0.028$ on ML1M with $n=4{,}832$ users;
$\rho=0.046$ on Steam with $n=31{,}836$ users). Utility-bucket means
moved in the same direction: higher-utility users were slightly less
negative, not more negative. We therefore keep the inversion claim at
the cross-dataset level and do not upgrade it to a user-level law.
% artifact: docs/reports/data/2026-07-04-followup09/per_user_utility_harm_report.md

\paragraph{Corrupted-bank gate response.}
We recomputed the frozen estimator on the existing frozen Beauty
token-dropout banks, using the same regenerated train split and the
same 4000-transition / seed-7 / 100-popularity-negative protocol as the
clean-bank artifact. Across the canonical corrupted ladder (30\%
through 50\% token dropout), $\Uds$ declined monotonically with dropout
rate within the corrupted banks, from $0.7133$ at the lightest
corrupted point to $0.7044$ at 50\% dropout, but every value remained
above the frozen $0.70$ shutoff; $\phids$ therefore stayed at $0$
throughout. Relative to the clean bank ($\Uds=0.7124$), the lightest
corrupted bank was slightly higher ($+0.0008$), so we record this as
observational evidence only rather than a strict monotone law. It does
not justify any frozen-parameter update, and on the archived Beauty
corruption banks the dataset-level gate never opens as text evidence
degrades.
% artifact: docs/reports/data/2026-07-04-followup10/beauty_corrupted_u_ds_phi_report.md

\paragraph{Pure out-of-sample check.}
ASO was frozen \emph{after} the four-point table above and evaluated
under a pre-registered rule fixed before its statistic existed: with
$\Uds{=}0.538$ and $\phids{=}1.0$, the frozen prediction was
$\Delta\mathrm{NDCG}@10 > 0$ against the ASO host. The observed result
was a miss: $\Delta\mathrm{NDCG}@10 = -0.0108$ on the test p2 selector.
We therefore record ASO as a pure out-of-sample prediction miss, not as
support for upgrading the gate claim.
% artifact: docs/reports/data/2026-07-02-gate0/aso_validation_report.md

\section{Experiments}
\label{sec:exp}

\paragraph{Setup.}
Four sequential-recommendation benchmarks (Steam, ML1M, Amazon Beauty,
and ATG (row-level regenerated split statistics in
Table~\ref{tab:dataset_stats}), rebuilt from raw data under a frozen protocol (raw
source, mapping, filtering, dedup, split rules recorded); all methods
share the regenerated splits and frozen banks. Metrics: HR@10 and
NDCG@10 under the production checkpoint selector; the host baseline is
the frozen PreferGrow run set. External baseline: DiffuRec, selected as
the single external closeout baseline and executed through a local
wrapper that reads the regenerated shared-protocol artifacts directly.
The dated E0 amendment confirms DiffuRec test-evaluator comparability
on all four datasets under the corrected common full-tail, row-weighted,
real-item test contract. This confirmation is limited to legacy-selected
frozen checkpoints under corrected test evaluation: validation
selection was not recomputed, and checkpoint-selection equivalence is
not claimed. Table~\ref{tab:validation} remains frozen in this cycle;
any later numeric replacement must move all systems together under the
same evaluator rather than selectively backfill one method~\cite{diffurec}.
% artifact: docs/reports/data/2026-07-10-evaluator-amendment/e0_evaluator_amendment.json
Hyperparameters ($\tau_0{=}0.2$, $g_{\max}{=}0.5$, $k{=}2.0$; frozen
$\phi$ per Eq.~\eqref{eq:phi}) and infrastructure in
Appendix~\ref{app:impl}.

\begin{table}[t]
\centering
\small
\begin{tabular}{lrrrr}
\toprule
Dataset & Items & Train rows & Val rows & Test rows \\
\midrule
Steam  & 9265  & 988517 & 81695 & 80651 \\
ML1M   & 3706  & 755782 & 98622 & 85405 \\
Beauty & 12101 & 17890  & 2236  & 2237  \\
ATG    & 11921 & 15529  & 1941  & 1942  \\
\bottomrule
\end{tabular}
\caption{Row-level regenerated split statistics used in the frozen protocol. Counts are distilled from the archived Gate0 component-summary artifact for the same paper-raw datasets.}
% artifact: docs/reports/data/2026-07-09-paper-dataset-stats-table.csv
\label{tab:dataset_stats}
\end{table}

\subsection{Frozen-gate validation across four datasets}
\label{sec:validation}

Under the frozen gate, $\phids$ fully opens Steam ($1.0$), barely opens
ATG ($0.117$), and closes Beauty and ML1M ($0$). Lemma~\ref{lem:tv}
bounds the proposal and one-step transition rows in total variation;
it does not bound end-to-end metric degradation. The pre-registered
predictions were fixed before launch: \emph{no-loss} on closed-gate
datasets, near-parity under the barely-open ATG gate, and preserved
gain on fully open Steam.

\begin{table}[t]
\centering
\small
\begin{tabular}{lccrl}
\toprule
Dataset ($\phids$) & Host & Ours & $\Delta$ & vs.\ prediction \\
\midrule
Steam (1.00)  & 0.0129 & 0.0149 & $+0.0020$ & sign \checkmark, mag.\ $\times$ \\
ATG (0.12)    & 0.0419 & 0.0305 & $-0.0114$ & barely-open parity $\times$ \\
Beauty (0.00) & 0.0333 & 0.0294 & $-0.0039$ & parity \checkmark \\
ML1M (0.00)   & 0.0910 & 0.0759 & $-0.0151$ & parity $\times$ \\
\bottomrule
\end{tabular}
\caption{Frozen-gate validation, test NDCG@10 (p2 selector), single
seed per the frozen protocol. Each dataset is read against its
pre-registered prediction; no uniform-superiority claim is made or
supported.}
% artifact: docs/reports/data/2026-07-06-gate1/sprint05_gate1_report.json
% (per-dataset core/current test p2 ndcg10 and deltas; manifests hash-checked)
\label{tab:validation}
\end{table}

The verdict is mixed and we report it as such. Fully open Steam kept a
positive sign but fell short of the reference magnitude; closed-gate
Beauty held its parity band. The closed-gate ML1M single-run miss
exceeds its measured host floor. The barely-open ATG single-run miss is
$4.37\times$ the provenance-limited observed spread across three
completed host runs; because all three frozen run manifests are absent,
that spread does not establish complete configuration parity or a fully
manifested noise floor. E0's corrected common test contract preserves
both miss decisions. Proposition~\ref{prop:reduction} is a statement about
the kernel objects at $g=0$; it does not by itself certify identity of
the complete training, selection, and evaluation implementation. We
therefore report one closed-gate implementation miss plus one
barely-open ATG miss as open implementation questions, and make no
claim that exact reduction has been demonstrated end-to-end. E1 instead
found the first production-path divergence at step 0: the canonical
core proposal parameter has different optimizer ownership in the host
and final-v2 paths. The downstream E2 attribution controls therefore
were not launched in this sprint. The theoretical proposition is
unaffected, and the pre-launch offline numeric suite matched $p_1$,
corruption rows, score-entropy terms, and reversal ratios at $g=0$.
% artifacts: docs/reports/data/2026-07-07-close02-ml1m-noise-floor/close02_ml1m_noise_floor_report.json
%            docs/reports/data/2026-07-10-evaluator-amendment/e0_evaluator_amendment.json
%            docs/reports/data/2026-07-10-e01-gzero-production-trace/e01_hard_stop.json
%            docs/reports/data/2026-07-10-close10-atg-noise-floor-provenance-limited/close10_atg_provenance_limited_report.json

\subsection{Controls and mechanism evidence}
\label{sec:mechanism}

The mechanism claims rest on the first-generation user-factor
instantiation (Appendix~\ref{app:gen1}), evaluated on Beauty under
token-dropout corruption training; all arms share splits, banks, and
selector.

\begin{table}[t]
\centering
\small
\begin{tabular}{lcc}
\toprule
Variant & val NDCG@10 & test NDCG@10 \\
\midrule
user-gated (aligned $u$) & 0.0235 & \textbf{0.0393} \\
anchor-only (no gating)  & \textbf{0.0238} & 0.0366 \\
$u$-shuffle              & 0.0212 & 0.0206 \\
global-$p$               & 0.0165 & 0.0289 \\
\bottomrule
\end{tabular}
\caption{Controls on Beauty (token-dropout regime, first-generation
instantiation; selector val NDCG@10 at $w{=}5$). Shuffling the
reliability signal across users collapses performance; a generic
global proposal bias does not reproduce the gain.}
% artifact: docs/reports/2026-06-28-beauty-token-dropout-control-results.md
\label{tab:controls}
\end{table}

The first-generation controls provide three bounded observations.
\emph{Aligned gating is
load-bearing}: permuting the reliability signal across users
($u$-shuffle) drops test NDCG@10 from $0.0393$ to $0.0206$; the signal
content, not the presence of a user-level dial, drives the effect.
\emph{Generic biases do not explain it}: a learned global proposal
without user conditioning reaches only $0.0289$. Under
description-blanking corruption, mean reliability drops by $0.0134$
with a corresponding rise in fallback mass; semantic token dropout has
the same direction ($-0.0031$), and the gated first-generation variant
degrades less than the ungated anchor from clean to corrupted test
conditions ($-0.0008$ vs.\ $-0.0043$ NDCG@10). These corruption
readouts do not validate a final-v2 corruption-response claim; the preregistered
E3 test was not launched after the E1 hard stop.
% artifacts: 2026-06-28-beauty-description-blank-corruption.md,
%            2026-06-28-beauty-token-dropout-end-to-end.md
Under the frozen two-factor gate, the clean-root v2 controls separate
closed-gate parity from opened-gate mechanism effects. On Beauty, where
$\phi(U_{ds})=0$, the \texttt{global\_p} arm is numerically identical
to the full run (both $0.0294$ test NDCG@10) and the $u$-shuffle arm
does not degrade it ($0.0330$), consistent with the dataset-level gate
being shut. On Steam, where $\phi(U_{ds})=1$, shuffling the user
reliability signal lowers test NDCG@10 by $0.0011$ relative to full,
while the \texttt{global\_p} arm stays close to the host
($+0.00035$ vs.\ core) and below the full run, so the deployed arm's
small positive gain is not explained by a dataset-agnostic bias alone.
The \texttt{text\_anchor\_only} arm is reported only as a high-variance
reference control rather than a replacement for the gated method; it
rises to $0.0302$ on Steam and sits slightly below full on Beauty.
% artifact: docs/reports/data/2026-07-06-sprint07/sprint07_control_table.csv

\pfig{First-generation corruption diagnostic: clean-vs-corrupted shift
of the gate and reliability, bucketed by clean reliability (regenerate
from archived diagnostics with the figure toolchain).}

\subsection{Injection-site analysis}
\label{sec:threeway}

We inject the identical signals at three sites: encoder (feature
fusion), loss (confidence weighting), and kernel (ours). In
first-generation runs the earlier injection sites were competitive in
clean regimes while kernel-side conditioning was strongest under
corruption; we report this as a regime observation scoped to that
instantiation, not a claim that the kernel site consistently wins. Per
Remark~\ref{rem:site}, the argument for the kernel site is not that it
wins every regime but that it is the only site where the fallback
guarantee is expressible.
% Scoped per spec 4.4; the planned three-way rerun under v2 was cancelled
% (2026-07-06 triage) and no cross-site superiority claim is made.

\section{From Gates to Schedules}
\label{sec:schedules}

Our gate acts as a \emph{constant} multiplier along the corruption
trajectory: the same $g$ applies at every noise level. The validation
outcomes above (an opened gate that keeps its sign but not its
reference magnitude, and a pure out-of-sample miss on ASO) are
consistent with this being too blunt an instrument: the value of
tilting corruption toward content plausibly differs between light
corruption (where the faded item is nearly recoverable) and heavy
corruption (where the proposal dominates the state distribution). A
simple algebraic fact makes the time-varying extension well-posed:
products of transitions of the form $\alpha I + (1-\alpha)\Pi_p$ with
\emph{different} proposals remain of the same form, with an effective
proposal that is an explicit convex combination of the per-step
proposals along the path. Closed-form marginals, score-entropy weights,
and the reduction property therefore survive schedule-dependent gating,
and the utility statistic $\Uds$ becomes the natural \emph{price} of a
corruption schedule rather than of a scalar gate. Developing this
evidence-priced schedule family, with the same pre-registration
discipline used here, is the direct continuation of this work.
% No empirical claim is made in this section (spec 5.1 boundary).

\section{Related Work}
\label{sec:related}

\paragraph{Diffusion recommenders.}
Continuous-diffusion recommenders inject Gaussian noise in a latent
space~\cite{dreamrec,diffurec,preferdiff}; discrete formulations
operate on the vocabulary directly~\cite{prefergrow,llada-rec}. Our
host is PreferGrow~\cite{prefergrow}, whose personalization is
score-side (condition dropout plus inference-time guidance against a
non-preference-user branch); its kernel is user-independent. We supply
the kernel-side half, anchored to their learned proposal.

\paragraph{Discrete diffusion LMs and adaptive corruption for
generative recommendation.}
Masked discrete diffusion language models~\cite{dllm-llada,sedd} now
back a fast-growing generative-recommendation line: parallel
semantic-ID generation~\cite{llada-rec}, confidence-ordered any-order
decoding~\cite{gr-diffgrm}, position-adaptive masking schedules
motivated by the finding that uniform masking is harmful
\cite{gr-admasked}, and emerging analysis~\cite{gr-mdtheory}. These
works adapt corruption to \emph{model-internal} signals (confidence,
position heuristics) and provide no external-evidence pricing and no
reduction guarantee when the adaptive signal is misleading. Our
results are complementary and cautionary: we give pre-registered
evidence that the value of evidence-shaped corruption can
\emph{invert} with evidence quality, a train-only statistic that
prices it before training, and a conditioning construction with exact
reduction at $g=0$ plus proposal and one-step transition-row TV bounds.
We view this as the reliability layer that
adaptive-corruption recommenders will need.

\paragraph{Shaped transition kernels in discrete diffusion.}
Structured corruption matrices date to D3PM~\cite{d3pm}; in
recommendation, CDRec~\cite{cdrec} shapes transitions with item
popularity and TDPM~\cite{tdpm} modulates token diffusion with temporal
preference structure. These are global or item-level signals; none
condition the kernel per user on evidence signals, none provide a
reduction guarantee, and none report the utility inversion: that
\emph{stronger} content evidence can make kernel shaping \emph{worse}.

\paragraph{Guidance and adaptivity.}
Classifier-free guidance~\cite{cfg} and its recommendation
uses~\cite{cfgrec} interpolate conditional and unconditional
\emph{scores} with a global weight; confidence-adaptive guidance exists
for language diffusion~\cite{acfg}. Our gate is the kernel-side
analogue with the weight set by measurable, history-only evidence
statistics; the interpolation target is the host's learned
proposal, which is what makes exact reduction and the stated
kernel-level TV bounds provable.

\paragraph{Reliability and side information in recommendation.}
Reliability-aware recommenders act at the representation, optimization,
or decision level~\cite{ugr,truthsr}; text-side sequential recommenders
fuse content into encoders~\cite{recformer,unisrec}. Neither line
conditions a diffusion transition law, and the hard-negative vs.\
false-negative distinction our gate encodes has no analogue at those
injection sites.

\paragraph{False negatives in negative sampling.}
That over-informative similarity turns sampled negatives into false
negatives is documented in contrastive representation
learning~\cite{debiased} and in negative sampling for implicit
collaborative filtering~\cite{srns}, where hard-negative miners must
avoid future positives. Our inversion transports this failure mode to
diffusion \emph{kernels}, where the ``negative sampler'' is the
transition law itself, entangled with the stationary distribution,
detailed balance, and the training objective
(Prop.~\ref{prop:collision}), and supplies a train-only statistic,
$\Uds$, that prices the risk before any training run.

\section{Conclusion and Limitations}
\label{sec:conclusion}

We personalized the transition kernel of discrete preference diffusion
with a two-factor, history-only evidence gate anchored to the host's
learned proposal. The formal fallback guarantee has a precise scope:
exact kernel reduction at $g=0$, proposal TV at most $g$, and one-step
forward transition-row TV at most $(1-\alpha_\sigma)g$. It is not an
end-to-end performance-downside guarantee. The design is built around
an empirical inversion we believe is of independent interest: because
kernel proposals are corruption laws, text evidence that predicts
positives is precisely the evidence a kernel must \emph{not} lean on.
$\Uds$ prices whether the kernel should admit the evidence before a
single training run.

\emph{Limitations.} The cross-dataset inversion rests on four points.
Under exchangeability, $1/24$ is only the descriptive fraction for one
specified order among $4!$ equally likely label orders, not a
confirmatory p-value. The within-dataset per-user follow-up does not
reproduce it: on ML1M and Steam, the utility-vs-%
$\Delta\mathrm{NDCG}@10$ association is weak and slightly positive
($\rho=0.028$ and $0.046$), so the inversion remains a dataset-level
observation rather than a user-level mechanism law. The frozen-gate
validation returned a mixed verdict: fully open Steam kept its sign but
not its reference magnitude, ASO missed outright, and the readout
contains one closed-gate miss (ML1M) plus one barely-open ATG
($\phids{=}0.117$) miss. The ML1M miss exceeds its measured host floor.
The ATG miss is $4.37\times$ a provenance-limited observed three-run
spread, but missing manifests preclude a complete configuration-parity
or fully manifested floor claim. E1 found a step-0 optimizer-ownership
divergence and consequently blocked E2 in this sprint. We do not claim
an end-to-end empirical demonstration of exact reduction or a stronger
causal attribution. The mechanism evidence comes from
the first-generation user-factor instantiation whose parameterization
differs from the final system (Appendix~\ref{app:gen1}). The $\phi$
thresholds, while pre-registered and robust across $65\%$ of a
post-hoc reporting-only sweep, remain a coarse hinge on limited
evidence; the ML1M bank enters with a single fixed-window history
length, so its null calibration reduces to one length bin; and the
guarantees are specific to kernels with a stationary
proposal; mask- or schedule-based discrete diffusions need their own
analysis (Sec.~\ref{sec:schedules}). An earlier metadata-sparsity
hypothesis was not supported and is reported as a negative result in
Appendix~\ref{app:results}.

% ======================= appendix =======================
\clearpage
\appendix

\section{Proofs}
\label{app:proofs}
Throughout, fix the user history $h$, the gate inputs $\ut$ and
$\Uds$, hence $g$ and the proposal $\pu$ of Eq.~\eqref{eq:pu}; $\pu$ is
strictly positive whenever $\pcore$ and (on its support) $\ch$ are.
Row-vector convention; $\alpha_\sigma=e^{-\sigma}$.

\begin{proposition}[Proposal-conditioned forward corruption]
\label{prop:forward-app}
Let $\Pi_u(i,j)=\pu(j)$ and $Q_u=\Pi_u-I$. Then
$P_{\sigma,u}=e^{\sigma Q_u}=\alpha_\sigma I+(1-\alpha_\sigma)\Pi_u$,
and for discrete steps with $\bar\alpha_t=\prod_{r\le t}\alpha_r$,
\[
  q_t(x_t{=}x\mid x_0,u)=\bar\alpha_t\,\mathbf{1}\{x{=}x_0\}
  +(1-\bar\alpha_t)\,\pu(x).
\]
\end{proposition}
\begin{proof}
$\Pi_u$ is an idempotent projection: $\Pi_u^2(i,j)=\sum_k
\pu(k)\pu(j)=\pu(j)=\Pi_u(i,j)$. Since $\Pi_u$ and $I$ commute,
$e^{\sigma Q_u}=e^{-\sigma}e^{\sigma\Pi_u}
=e^{-\sigma}\big(I+(e^{\sigma}-1)\Pi_u\big)
=\alpha_\sigma I+(1-\alpha_\sigma)\Pi_u$.
Products telescope by idempotence:
$\prod_t\big(\alpha_t I+(1-\alpha_t)\Pi_u\big)
=\bar\alpha_t I+(1-\bar\alpha_t)\Pi_u$; take row $x_0$, column $x$.
\end{proof}

\begin{proposition}[Conditional Markovianity]
\label{prop:markov-app}
Conditional on the fixed gate inputs, $(x_t)$ generated by
$P_{\Delta\sigma_t,u}$ is a (time-inhomogeneous, per-step-size) Markov
chain; under a fixed step size it is generated by the single generator
$Q_u$, hence a homogeneous Markov semigroup.
\end{proposition}
\begin{proof}
With the gate inputs fixed, $\pu$ and $Q_u$ are deterministic; each
transition reads only $(x_t,\Delta\sigma_t)$, never earlier states, so
the conditional law of $x_{t+1}$ given the past equals
$P_{\Delta\sigma_t,u}(x_t,\cdot)$. Without conditioning, the mixture
over users need not be a single Markov chain.
\end{proof}

\begin{proposition}[Stationary and limiting proposal]
\label{prop:stationary-app}
$\pu P_{\sigma,u}=\pu$ for all $\sigma\ge0$, and for any initial law
$\mu$, $\mu P_{\sigma,u}=\alpha_\sigma\mu+(1-\alpha_\sigma)\pu\to\pu$
as $\sigma\to\infty$. For finite $\sigma$ the marginal is the stated
mixture, not $\pu$ itself.
\end{proposition}
\begin{proof}
$\mu\Pi_u=\pu$ for any distribution $\mu$; substitute into
$P_{\sigma,u}=\alpha_\sigma I+(1-\alpha_\sigma)\Pi_u$.
\end{proof}

\begin{proposition}[Detailed balance]
\label{prop:balance-app}
$\pu(i)P_{\sigma,u}(i,j)=\pu(j)P_{\sigma,u}(j,i)$ for all $i,j$.
\end{proposition}
\begin{proof}
Diagonal terms are trivial; for $i\ne j$ both sides equal
$(1-\alpha_\sigma)\pu(i)\pu(j)$.
\end{proof}

\begin{proposition}[Analytic reverse ratio replaces only the proposal]
\label{prop:reverse-app}
For $\delta<\infty$, $P_{\delta,u}$ is invertible with
$P_{\delta,u}^{-1}=\alpha_\delta^{-1}I+(1-\alpha_\delta^{-1})\Pi_u$;
hence for any row vector $r$,
$rP_{\delta,u}^{-1}=e^{\delta}r+(1-e^{\delta})\langle
r,\mathbf{1}\rangle\,\pu$. The host's stationary-reversal operator is
recovered with $\pu$ in place of the global proposal; exact
non-stationary reverse kernels additionally depend on the current
marginal and remain the score model's job.
\end{proposition}
\begin{proof}
Multiply the two matrices and use $\Pi_u^2=\Pi_u$:
$[\alpha I+(1-\alpha)\Pi][\alpha^{-1}I+(1-\alpha^{-1})\Pi]=I$. The
row-vector form follows from $r\Pi_u=\langle r,\mathbf{1}\rangle\pu$.
\end{proof}

\begin{proposition}[Exact reduction; Prop.~\ref{prop:reduction} in the
main text]
\label{prop:reduction-app}
If $g=0$ then $\pu=\pcore(\star)e_\star
+(1-\pcore(\star))\pcorer=\pcore$, hence $\Pi_u=\Pi_{\pcore}$,
$Q_u=Q_{\pcore}$, and every object in
Props.~\ref{prop:forward-app}--\ref{prop:reverse-app}, as well as the
score-entropy weights (which depend on the kernel only through $\pu$
and the schedule), coincides with the host's.
\end{proposition}
\begin{proof}
Immediate: $\pcorer$ renormalizes $\pcore$ off the pseudo-item, so
recombining with mass $\pcore(\star)$ restores $\pcore$ exactly; all
listed objects are functions of $(\pu,\text{schedule})$ alone.
\end{proof}

\begin{lemma}[TV bound; Lemma~\ref{lem:tv} in the main text]
\label{lem:tv-app}
$\TV(\pu,\pcore)\le g$ and
$\TV\big(P_{\sigma,\pu}(i,\cdot),P_{\sigma,\pcore}(i,\cdot)\big)
\le(1-\alpha_\sigma)\,g$.
\end{lemma}
\begin{proof}
$\pu-\pcore=(1-\pcore(\star))\,g\,(\ch-\pcorer)$ on real coordinates
and $0$ on $\star$, so
$\TV(\pu,\pcore)=\tfrac12\|\pu-\pcore\|_1
=(1-\pcore(\star))\,g\,\TV(\ch,\pcorer)\le g$.
Rows of the transitions differ by $(1-\alpha_\sigma)(\pu-\pcore)$.
\end{proof}

\begin{proposition}[Target collision; Prop.~\ref{prop:collision} in
the main text]
\label{prop:collision-app}
$\Pr[x_t{=}y\mid x_0{\ne}y]=(1-\bar\alpha_t)\pu(y)$, and $\pu(y)$ is
affine in $g$ with slope
$(1-\pcore(\star))(\ch(y)-\pcorer(y))$.
\end{proposition}
\begin{proof}
By Prop.~\ref{prop:forward-app},
$q_t(x_t{=}y\mid x_0,u)=\bar\alpha_t\mathbf{1}\{y{=}x_0\}
+(1-\bar\alpha_t)\pu(y)$; with $x_0\ne y$ the indicator vanishes.
Substituting Eq.~\eqref{eq:pu} at coordinate $y\ne\star$ gives
$\pu(y)=(1-\pcore(\star))[(1-g)\pcorer(y)+g\,\ch(y)]$, which is affine
in $g$ with the stated slope; it is increasing iff
$\ch(y)>\pcorer(y)$.
\end{proof}

\section{Scope and boundary conditions}
\label{app:boundary}
The guarantees require: (i) gate inputs computed from history only and
held fixed along each trajectory; (ii) self-resampling corruption of
the form \eqref{eq:forward}; (iii) the pseudo-item coordinate, when
present, carried by every kernel-level object with its mass taken from
$\pcore$ and never modified by the gate; (iv) forced replacement,
top-$k$ truncation, or state-dependent masking break detailed balance;
(v) hard-masked items restrict irreducibility claims to the support;
(vi) the reverse-ratio statement concerns the stationary-reversal
analytic operator, not the exact non-stationary reverse kernel; and
(vii) if the gate leaks batch labels at training time, the certified
object is a different, data-dependent kernel.

\section{The $\Uds$ estimator and pre-registration}
\label{app:uds}
Frozen estimator: per dataset, sample $4{,}000$ train transitions
(seed 7); the history vector is the normalized mean of valid item
embeddings from the production bank; $\Uds$ is the mean probability
that anchor similarity ranks the true next item above each of $100$
popularity-sampled negatives (train next-item frequencies);
uniform-negative and coherence-quartile variants are reported as
diagnostics. Artifacts record bank and split hashes. The hinge
parameters $(0.70, 0.10)$ of Eq.~\eqref{eq:phi} were frozen from a
legacy-bank precheck before production statistics existed and were not
adjusted afterwards.
% artifact: docs/reports/data/2026-07-02-gate0/gate0_text_utility_report.json

\section{First-generation instantiation}
\label{app:gen1}
The first-generation runs (Table~\ref{tab:inversion}, right column;
Table~\ref{tab:controls}) used the user-factor-only gate with the
pre-revision parameterization (three-dial adaptive mapping) and no
dataset factor. They are mechanism evidence and failure-mode
documentation, not benchmarks of the final system. The richer user
factor combined behavior-text coherence, text completeness, and
history-stability proxies, then mapped them to three dials (anchor
temperature, popularity mixing, and pseudo-item mass) without the
dataset utility factor and without anchoring to the host's learned
proposal. Its ML1M failure is precisely what the final design removes:
at the typical reliability level its pseudo-item mass departed far from
the host's learned value, and no dataset-level factor prevented content
tilt on the highest-utility bank. The final system replaces the three
dials with the single gate $g$, freezes the pseudo-item mass to the
host's, and adds $\phids$. First-generation numbers are therefore
evidence about aligned user gating, limited corruption responses, and
failure modes, not evidence of a final-v2 corruption-response claim or a benchmark
of the final system.

\section{Implementation and infrastructure}
\label{app:impl}
Hyperparameters: $\tau_0{=}0.2$, $g_{\max}{=}0.5$, $k{=}2.0$, frozen
$\phi(0.70,0.10)$; selector: validation NDCG@10 (p5); banks:
sentence-t5-xl, frozen and hash-recorded; splits regenerated from raw
data under a recorded protocol; all validation runs carry manifests
recording config, seed, bank/null-curve/split/$\Uds$-artifact hashes.
Model selection used validation only; test metrics were logged during
development. Accordingly, we do not describe the test split as an
untouched final holdout.
Code and data for reproducibility accompany the submission as
supplementary material per AAAI-27 policy; the reproducibility
checklist is submitted separately. A fallback compile validation has
already been exercised with Tectonic~0.16.9 on the remote clean-root
closeout environment, yielding an article-mode PDF without undefined
citations; camera-style page-count verification still awaits the
AAAI-27 author kit, and the supplementary bundle will package the dated
artifacts, manifests, and checklist required by the policy.
% artifact: issues/2026-07-06_evidence-priced-schedule-and-closeout.csv (CLOSE-07 notes: compile_validation_first_pass, compile_result, compile_caveat)
% artifact: docs/reports/2026-07-07-aaai27-handoff-playbook.md

\section{Negative results}
\label{app:results}
The metadata-sparsity robustness hypothesis was not supported and is
reported here as a negative result rather than omitted. The original
claim was that completeness bucketing alone might predict per-user
gains; the archived results did not support that reading, so the main
text no longer uses metadata sparsity as positive evidence and instead
treats it as a negative appendix-side finding.

\bibliography{references}
\IfFileExists{aaai27.bst}{\bibliographystyle{aaai27}}{\bibliographystyle{plain}}

\end{document}
